\documentclass[11pt]{article}
\usepackage{amsmath,amssymb}
\title{Room 18 (seat: direct) --- accumulation at $N=6$ is already proved; what remains open is
strictly narrower than the contour construction}
\date{2026-09-27}
\begin{document}
\maketitle

\section*{Verdict up front}

\textbf{GENUINE PROGRESS TOWARD ACCUMULATION --- in fact the accumulation statement itself is
already closed, unconditionally for $G^T_0$ and conditionally (on \texttt{thm:model},
\texttt{thm:RJ}, exactly as everywhere else in the $U,V$ theory) for $U,V$.} This is not a new
theorem: it is already written, as a proof, inside \texttt{paper2a.tex} (Theorem~\texttt{thm:arc}
and Corollary~\texttt{cor:arcUV}), and it was already flagged in two places in this repository
(\texttt{P1i\_lemma.tex} lines 27/186, and the older
\texttt{P1\_N6/room1\_seat\_erdos.tex}). What Room~18 contributes is: (1) tracing the exact
mechanism and confirming it really does cover the $N=6$ point with no gap, (2) noting explicitly
that the entire $N=6$-v2 track (Rooms 8, 11, 13, 14, 16, 17 in this directory) never connected to
this fact and so kept treating ``accumulation at $N=6$'' as blocked on the missing
\texttt{prop:qi}-style contour, and (3) precisely re-scoping what actually remains open: not
accumulation, but the strictly stronger \emph{direct tie-ray statement at $N=6$ itself}. A
secondary finding (\S\ref{sec:secondary}) reassesses Room~11's stated wall against Room~17's fix
and finds Room~11 overstated it: Room~17's bound does remove the obstruction on the bulk of the
constant-height legs ($\operatorname{Re}x>0$), not just at isolated points, and the real remaining
gap is narrower still.

\section*{1. The closure mechanism, checked directly against paper2a.tex}

\texttt{paper2a.tex}, \S\ref{sec:arc} (``The middle arc''), Theorem~\texttt{thm:arc}:
\begin{quote}
The zeros of $1-\Sigma_1=\cos(Z;q^2)$ in $0<|q|<1$, which are the poles of $G^T_0$, accumulate at
\emph{every point} of the arc $\{|q|=1,\ \arg q\in[0.286\pi,1.714\pi]\}$.
\end{quote}
The proof sketch (checked in full, paper2a.tex lines 2508--2534) is: for every $\zeta=e^{2\pi ia/N}$
with $\gcd(a,N)=1$, $\frac aN\in[0.143,0.857]$, and $N\ge17$ (or $N\ge16$ if $\frac aN\in[\frac16,
\frac56]$), the \texttt{prop:qi}-style construction (P1f/P1h's Theorems M$'$, B$'$, landscape Lemma
LS$(\ell_1)$, Rouch\'e against $1+\varrho e^{\Delta V/t}$) is actually carried out and gives an
explicit sequence of zeros $q^\ast_j\to\zeta$. This makes each such $\zeta$ itself a genuine
accumulation point of the actual pole set $P$ of $G^T_0$. The proof's last sentence is the step that
matters here: \emph{``These $\zeta$ are dense in the arc, and the accumulation set is closed.''}

This is a standard, elementary closure argument and it is fully general (it needs nothing about
$N=6$ specifically):
\begin{itemize}
\item The derived set $P'$ (set of accumulation points) of any subset $P\subset\mathbb C$ is closed.
(Standard: if $p$ is a limit of points of $P'$, each of which is itself a limit of points of $P$,
a diagonal argument produces points of $P$ arbitrarily close to $p$.)
\item Let $S=\{\zeta=e^{2\pi ia/N}:\gcd(a,N)=1,\ N\ge17\ (\text{or }N\ge16\text{ on }[\tfrac16,
\tfrac56]),\ \tfrac aN\in[0.143,0.857]\}$. The theorem's proof shows $S\subseteq P'$.
\item $S$ is dense in the arc $\{\arg q\in[0.286\pi,1.714\pi]\}$: for any target angle
$2\pi r$ with $r\in[0.143,0.857]$ and any $\epsilon>0$, reduced fractions $a/N$ with $N$ as large as
one likes and $|a/N-r|<\epsilon/(2\pi)$ exist (elementary density of $\mathbb Q$), and for $N$ large
enough these automatically satisfy $N\ge17$ (or the sharper $N\ge16$ threshold if $r\in[\frac16,
\frac56]$).
\item Since $P'$ is closed and contains the dense set $S$, $P'$ contains the closure of $S$, which
is the whole arc.
\end{itemize}
\textbf{Hence every point of the arc, including $\zeta_0=e^{2\pi i/6}=e^{i\pi/3}$ ($N=6$, $a=1$)
itself, is an accumulation point of the poles of $G^T_0$.} The point $a/N=\frac16=0.1\overline6$
satisfies $0.143<0.1\overline6<0.857$ strictly (it is well inside the arc, not a boundary case for
this theorem), so there is no edge-case subtlety here at all: it is exactly the kind of point the
theorem's closure step was built to reach.

\section*{2. The same argument transfers to $U,V$}

\texttt{Corollary~cor:arcUV} (paper2a.tex lines 2538--2564), conditional on \texttt{thm:model} and
\texttt{thm:RJ} exactly as every other $U,V$-pole statement in the paper, states: for every
$\zeta=e^{2\pi ia/N}$ with $\gcd(a,N)=1$, $N\ge16$, $\frac16\le\frac aN\le\frac56$, the poles of $U$
and $V$ accumulate at $\pm\sqrt\zeta,\pm\sqrt{\bar\zeta}$, and \emph{``the points $\pm e^{\pm i\pi a/N}$
are dense in the two arcs and the accumulation set is closed''} --- the identical mechanism,
applied this time to the arc $\arg x\in[\frac\pi6,\frac{5\pi}6]\cup[\frac{7\pi}6,\frac{11\pi}6]$.
Here $a/N=\frac16$ is the closed left endpoint of $[\frac16,\frac56]$, not an interior point, but
this does not affect the argument: density of $S$ up to and including approach from inside the
interval (e.g.\ $a_k/N_k\downarrow\frac16$ with $N_k\ge16$, $a_k/N_k\ne\frac16$) still makes
$\frac16$ a limit point of $S$, and $P'$ closed still forces $\frac16\in$ the covered set.
\textbf{So the poles of $U$ and $V$ accumulate at $x=e^{i\pi/6}$ (and $\pm e^{\pm i\pi/6}$'s
partners $e^{5i\pi/6},e^{-i\pi/6},e^{-5i\pi/6}$), which are exactly $\pm\sqrt{\zeta_0}$,
$\pm\sqrt{\bar\zeta_0}$ at $\zeta_0=e^{i\pi/3}$: the $N=6$ point.} This is conditional exactly as
stated in \texttt{rem:arcUV} (on \texttt{thm:model}, \texttt{thm:RJ}), no more and no less than
every other point on the two arcs already claimed as proved in the paper.

\section*{3. Why this was missed in this track}

\texttt{P1i\_lemma.tex} already states this (line 27: ``The bare statement `the zeros of $B$
accumulate at $\zeta$' is \emph{already} a consequence of P1h... this includes $\frac16$ and
$\frac56$''; line 186: ``Accumulation of the zeros of $B$ at $e(\frac16)$, and of the poles of
$U,V$ at $\pm\sqrt{e(\pm\frac16)}$, holds anyway by closure from P1h... What is open at $\frac16$ is
the direct tie-ray statement''), and the older \texttt{P1\_N6/room1\_seat\_erdos.tex} restates it
almost verbatim as its own calibration point before going on to a different (elementary,
Rouch\'e-on-a-small-disk) task. But none of Rooms 8, 11, 13, 14, 16, 17 in \emph{this} directory
(\texttt{P1\_N6\_v2}) cite \texttt{P1i\_lemma.tex} or this closure fact anywhere (checked by
\texttt{grep -l "closure\textbackslash|P1i\_lemma\textbackslash|P1h" *.tex}: only Room~8's files
match, and only for an unrelated ``small closure'' of two certified zeros, explicitly disclaimed
there as \emph{``not a density/accumulation theorem''}). Room~11's and Room~17's own framing (``What
this does and does not achieve: NOT achieved, still OPEN: $N=6$ itself'') is accurate about the
\emph{contour construction}, but conflates that with accumulation, which is a strictly weaker
consequence already available for free.

\section*{4. What is actually still open}

The genuinely open target, precisely: an explicit, self-contained, \texttt{prop:qi}-style statement
\emph{at $N=6$ itself} (not reached via a limit over other $N$): a sequence $t^\ast_j\to0$ with
$\arg t^\ast_j\to\theta^\ast=0$, $q^\ast_j=\zeta_0e^{-t^\ast_j}$ the \emph{unique} zero of
$1-\Sigma_1$ (or of $B$, in the P1f/P1i normalization) in an explicit disc $D_j$ shrinking like
$1/j$, with the two-saddle asymptotic $t^\ast_j=t^0_j+O(j^{-3})$ and explicit amplitudes, exactly as
Corollary~\texttt{cor:Ui} gives at $N=4$. This is strictly stronger than accumulation (it pins down
\emph{which} points converge to $\zeta_0$, along which ray, at which rate, self-containedly at
$N=6$, rather than merely that $\zeta_0$ is a limit of points where such statements hold for
\emph{other}, larger $N$). It is this narrower target that Rooms 11--17 correctly found blocked, and
it remains blocked here too: this room does not build the $N=6$ analogue of
$\Gamma_\pm^{(6)},\Pi_0^{(6)}$, the Part-(ii) crossing-region estimate, or the six-class landscape
certificate. \textbf{Nothing in \S\S1--3 above closes that.}

\section*{5. Secondary finding: Room 17's fix reaches further than Room 17/Room 11 credited}
\label{sec:secondary}

Room~11 (\texttt{room11\_seat\_erdos.tex}, Step~2--3) argued the direct contour construction stalls
because the two constant-height legs $\Gamma_-^{(6)}$ (through $x_{-3}$, $\operatorname{Im}x=
-\pi/6$) and $\Pi_0^{(6)}$ (through $x_{-1}$, $\operatorname{Im}x=\pi/6$) sit exactly on two of the
five $N=6$ resonance loci for their \emph{entire length}, and concluded: ``the ray lemma's proof
gives no lower bound on the separation constant $\delta$ \emph{anywhere along these legs}.'' Tracing
through the classical derivation ($\delta\ge\mathcal D(a',\varphi)$, $\mathcal
D(a',\varphi)=\inf_{\rho\ge0}\max\{\sigma(\varphi-\rho\sin\theta),Y(\rho)\}$,
$Y(\rho)=1-e^{-a'-\rho\cos\theta}$) at the resonant class ($\varphi=0$) with $\theta=\theta^\ast=0$
(so $\sin\theta=0$ identically): the rotation term $\sigma(\varphi-\rho\sin\theta)=\sigma(0)=0$ for
\emph{every} $\rho$ (the resonance really does kill it completely, confirming Room~11's diagnosis of
\emph{why} the classical angular margin is useless here), but $\mathcal D(a',0)=\inf_\rho\max\{0,
Y(\rho)\}=\inf_\rho Y(\rho)$. Since $Y$ is non-decreasing in $\rho$ whenever $\cos\theta\ge0$
(exactly Room~17's monotonicity observation, independently re-derived here), this infimum is
$Y(0)=1-e^{-a'}$, which is \emph{strictly positive whenever $a'=2\operatorname{Re}x>0$}, and $\to0$
only in the genuine limit $\operatorname{Re}x\to0$. So the correct statement is: the separation
constant is positive and explicit ($\ge1-e^{-2\operatorname{Re}x}$) at \emph{every} point of both
constant-height legs with $\operatorname{Re}x>0$ --- which is all of both legs except the single
point where each one crosses $\operatorname{Re}x=0$ on its way out from $x_c$ (which itself has
$\operatorname{Re}x_c<0$). This is not ``no lower bound anywhere along these legs''; it is a clean,
uniform, positive lower bound on the entire $\operatorname{Re}x>0$ portion, including both bounded
saddle windows (Room~17's own $N=6$-saddle number, $d(x_{-1})\ge0.5025\ldots$, is the special case
$\operatorname{Re}x=\operatorname{Re}(x_{-1})=0.349119\ldots$ of exactly this general statement, not
an isolated coincidence).

\textbf{What this narrows the remaining gap to.} The genuine obstruction is therefore not ``the ray
lemma fails throughout the resonant legs'' but specifically the single transition region
$\operatorname{Re}x\lesssim0$ where each leg passes from the start box near $x_c$ through
$\operatorname{Re}x=0$ into the region just analysed. In the $N=4$ template this is exactly the
region handled by \texttt{lem:qi-bounds}(ii) --- the Poisson-summation/triple-product argument
($\vartheta(z)=(z;P)_\infty(P/z;P)_\infty$, giving $\Upsilon(x)/r+M'(x)+\ldots$), which is a
genuinely different piece of machinery from the ray lemma and was \emph{not} touched by Rooms
11--17's work on $d(x)$, $d'(x)$. Building the $N$-generic ($N=6$, six residue classes) analogue of
part (ii) is therefore the actual remaining obstruction for the direct contour, not the separation
constant on the bulk of the legs (which this room now considers resolved, modulo writing out
part~(i)'s other, non-$\delta$ error constants $c_\ell,K_n$ for $N=6$, not attempted here).

\section*{Honest rigor labels}

\begin{itemize}
\item \textbf{PROVED (already, in paper2a.tex, unconditionally):} the poles of $G^T_0$ accumulate at
$\zeta_0=e^{i\pi/3}$ ($N=6$). Mechanism: closure of the dense accumulation-point set $S$ established
by Theorem~\texttt{thm:arc}, applied at $a/N=\frac16\in(0.143,0.857)$. No new computation was needed;
this room verified the argument line-by-line against paper2a.tex and confirmed $\frac16$ is strictly
interior to the covered arc.
\item \textbf{PROVED (already, in paper2a.tex, conditional on \texttt{thm:model},\texttt{thm:RJ},
same conditionality as the rest of the $U,V$ theory):} the poles of $U,V$ accumulate at $\pm\sqrt{
\zeta_0},\pm\sqrt{\bar\zeta_0}$. Mechanism: the identical closure argument via
Corollary~\texttt{cor:arcUV}, with $\frac16$ as the (still-covered, by density-from-inside) left
endpoint of $[\frac16,\frac56]$.
\item \textbf{OPEN, correctly identified as such by Rooms 11/13/14/16/17, and still open after this
room:} the direct, self-contained $N=6$ tie-ray statement (the literal \texttt{prop:qi} analogue at
$N=6$ itself, with an explicit two-saddle asymptotic for $t^\ast_j$).
\item \textbf{HEURISTIC/clarifying, this room, \S\ref{sec:secondary}:} Room~11's ``no lower bound
anywhere along the legs'' is an overstatement; Room~17's bound (independently re-derived here via
the classical $\mathcal D(a',\varphi)$ formula rather than asserted from scratch) gives a clean
positive separation constant on all of $\operatorname{Re}x>0$ on both constant-height legs, not just
at the saddle points. This narrows, but does not close, the remaining construction: what is missing
is specifically an $N=6$/six-class analogue of \texttt{lem:qi-bounds}(ii) for the
$\operatorname{Re}x\lesssim0$ crossing region, not the ray lemma itself.
\end{itemize}

\section*{Why routes (a)/(b) of the brief were not pursued further here}

Given \S\S1--4, the brief's stated goal (``giving a sequence of zeros $t^\ast_j\to$ the saddle,
hence accumulation at $e(1/6)$'') is achieved by a two-line closure argument already sitting in the
paper, without needing route~(a)'s full contour or route~(b)'s longer certified sequence. Extending
Room~1's direct Rouch\'e/winding-number certification (route~(b)) to more zeros $j=1,\ldots,J$ would
still be worthwhile \emph{for the direct tie-ray statement} (\S4), since it is independent evidence
for that narrower, still-open claim, but it would not add anything to the accumulation question
itself, which this room found already closed. No new heavy computation was run in this room; the
finding is a proof-reading/cross-referencing result, checked against the primary sources
(\texttt{paper2a.tex} lines 2496--2564, \texttt{P1i\_lemma.tex} lines 14--31 and 178--199) rather
than against any room summary.

\end{document}
